\documentstyle[%


righttag%


,11pt%


,amssymb%


,russian%               remove in Eng.


,izvru%                 Set izven% in Eng.


]{amsart}





% matevossian@cswebmail.com; filimi@yandex.ru





\newcommand{\dd}{\partial}


\newcommand{\bbR}{\Bbb R}


\newcommand{\Rn}{\bbR^n}


\newcommand{\eps}{\varepsilon}


\newcommand{\grad}{\bigtriangledown}


\newcommand{\dist}{\operatorname{dist}}


\newcommand{\sgn}{\operatorname{sgn}}


\newcommand{\loc}{\operatorname{loc}}


\newcommand{\Vol}{\operatorname{Vol}}


\newcommand{\const}{\operatorname{const}}


\newcommand{\stdsum}{\sum\limits_{i,j=1}^{n}}


\newcommand{\mgrad}{|\grad w_\eps(x,t)|}


\newcommand{\tgrad}{|\grad\theta_\eps(x)|}


\newcommand{\tmd}[1]{\theta_\eps^{#1}(x)}


\newcommand{\tm}{\tmd{}}


\newcommand{\tmp}{\tmd{p}}


\newcommand{\tts}[1]{}





%\hoffset=-15mm%                  For Eng.


\setcounter{page}{29}


\god{2005}


\nomer{9~(520)} %                        Remove in Eng.


%\nomer{Vol.\,49, No.\,9}%               Set in Eng.


%\str{\pageref{begin}--\pageref{end}}%  Set in Eng.


\udk{517.95}





\author{\it О.А.\,Матевосян, И.В.\,Филимонова}


% NB:   ^^ remove in Eng.


\title{Об усреднении слабо-нелинейных параболических операторов в


перфорированном цилиндре}


\thanks{}





\date{}


%\pagestyle{myheadings}%         Set in Eng.


\pagestyle{plain} %              Remove in Eng.


\begin{document}


%\thispagestyle{plain}%          Set in Eng.


\maketitle


\label{begin}








Изучению вопросов теории усреднения для перфорированных областей


посвящено много работ (см.~\cite{BLP}--\cite{MKh}), в которых на


поведение решений при приближении к дырам ставятся какие-либо


ограничения, например, граничные условия. Однако для усреднения


решений уравнений с достаточно сильным нелинейным поглощением


такие ограничения не нужны. Подобная ситуация была обнаружена и в


теории устранения особенностей. В связи с этим отметим работы


\cite{BF}--\cite{IOK}. Вопросы усреднения полулинейных


эллиптических задач в перфорированных областях были изучены в


работах \cite{MP1}, \cite{MP2} с использованием оценки


Кондратьева--Ландиса (\cite{KL}, \S\,1, п.\,3).





В данной работе для слабо-нелинейного дивергентного


параболического уравнения второго порядка с младшим членом,


растущим по неизвестной функции степенным образом, доказывается,


что последовательность решений в перфорированном цилиндре


стремится к решению в неперфорированном цилиндре, если радиусы


выброшенных цилиндров стремятся к нулю со скоростью, зависящей от


показателя степени в младшем члене. Доказательства теорем


данной работы основываются на аналогичной оценке решения


полулинейного парабoлического уравнения через расстояние от точки


до границы области в соответствующей метрике, полученной в работе


\cite{Ch}.





Будем пользоваться следующими обозначениями:


\begin{gather*}


x=(x_1,\dotsc,x_n)\in\Bbb{R}^{n},\ \ n>2,\ \ u=(u_1,\dotsc,u_n),


\ \ u_x=(u_{x_1},\dotsc,u_{x_n}),\\


|u_x|^2=\sum\limits_{i=1}^n(u_{x_i})^2,


\ \ u_t=\dd u/\dd t, \ \ 0\leq t<T<\infty;


\end{gather*}





$\Omega \subset \Rn$ --- произвольное открытое подмножество $\Rn$


с границей $\dd\Omega$,


$\overline\Omega=\Omega\bigcup\dd\Omega$ --- замыкание $\Omega$;


$Q= \Omega\times(0,T)$ --- цилиндр в пространстве точек


$(x,t)\in\Bbb{R}^{n+1}$, $x\in\Omega$, $t\in[0,T]$,


$0<T<\infty$;





$S(Q)=\partial\Omega\times[0,T]$ --- боковая поверхность цилиндра


$Q\subset\Bbb{R}^{n+1}$;


$\Gamma(Q)=S(Q)\cup\Upsilon_0(\Omega)$ --- параболическая граница


цилиндра $Q\subset\Bbb{R}^{n+1}$, где


$\Upsilon_t(\Omega)=\overline{\Omega}\times\{t\}$ \; $\forall t\in


[0,T]$;





$Q(x_0,t_0,R,T)=B(x_0,R)\times[t_0,t_0+T]$ --- цилиндр в


$\Bbb{R}^{n+1}$ радиуса $R$ и высоты $T$, где $B(x_0, R)=


\{x\in \Rn : |x-x_0| \le R\}$ --- замкнутый шар в $\Rn$ радиуса


$R$ с центром в точке $x_0$;





$\cal{C}(Q)$ --- совокупность непрерывных в $Q$


функций;





$L_{p,r}(Q)$ --- пространство, состоящее из всех измеримых на $Q$


функций с конечной нормой


$$


\|u;L_{p,r}(Q)\|=\bigg(\int_{0}^{T}\bigg(\int_{\Omega}


|u(x,t)|^{p}dx\bigg)^{r/p}dt\bigg)^{1/r},\ \ p\geq 1,\ \ r\geq 1;


$$


$L_{p,r}(Q)$ будем обозначать через $L_{p}(Q)$ и норму


$\|\,\cdot\,;L_{p,r}(Q)\|$ через $\|\,\cdot\,;L_{p}(Q)\|$ при $p=r$; 


$W^{1, 1}_p(Q)$ --- пространство Соболева, состоящее из элементов


$L_p(Q)$, имеющих обобщенные производные по $x$ и по $t$ с нормой


$\|u;W^{1, 1}_p(Q)\|=\|u; L_{p}(Q)\|+\|u_x;L_{p}(Q)\|+\|u_t;L_{p}(Q)\|;$





$W^{1, 0}_p(Q)$ --- пространство Соболева, состоящее из элементов


$L_p(Q)$, имеющих обобщенные производные по $x$ с нормой


$\|u;W^{1, 0}_p(Q)\|=\|u; L_{p}(Q)\|+\|u_x; L_{p}(Q)\|$.





Пусть $L$, $\cal{L}$ --- операторы вида


$$


L \equiv \stdsum\frac{\dd}{\dd x_i}\bigg(a_{ij}(x)


\frac{\dd}{\dd x_j}\bigg), \quad


\cal{L}\equiv L-\frac{\dd }{\dd t},


$$


где $(x,t)=(x_1,\dotsc,x_n,t)\in\Bbb{R}^{n+1}$, коэффициенты


$a_{ij}(x):\Bbb{R}^n\rightarrow\Bbb{R}^1$ --- ограниченные


измеримые функции такие, что $a_{ij}(x)\equiv a_{ji}(x)$ для всех


$x\in\Bbb{R}^n$, и выполняются условия


\begin{equation}


\lambda^{-1}|\xi|^2\le\stdsum a_{ij}(x) \xi_i \xi_j \le


\lambda|\xi|^2, \quad \lambda=\const>0,


\label{lmbd}


\end{equation}


для всех $\xi=(\xi_1,\dotsc,\xi_n)\in\Rn$,


$|\xi|^2=\sum\limits_{i=1}^n\xi_i^2$.





Будем изучать решение задачи


\begin{alignat}{2}


\cal{L}u(x,t)&=f(x,t,u(x,t)), &\quad (x,t)&\in Q=\Omega\times(0,T), \label{P}\\


u(x,t)&=\varphi(x,t), &\quad (x,t)&\in \Gamma(Q)=S(Q)\cup\Upsilon_0(\Omega),


\label{B}


\end{alignat}


где $\varphi(x,t)\in W_{2}^{1,1}(Q)$, а функция


$f(x,t,u):\Bbb{R}^{n+1}\times\Bbb{R}^1\rightarrow\Bbb{R}^1$


предполагается измеримой и такой, что композиция $f(x,t,u(x,t))$


является измеримой локально ограниченной функцией переменных


$(x,t)$ при любой функции $u(x,t)\in W_{2}^{1,1}(Q)\cap C(Q)$.








\begin{defn}


Функция $u(x,t)\in W^{1,1}_{2,\loc}(Q)\bigcap C(Q)$ называется


{\it обобщенным решением задачи} \eqref{P}, \eqref{B},


если для любой функции $\psi(x,t)\in\overset{0}{W}{}_{2}^{1,0}(Q)$


она удовлетворяет интегральному соотношению


\begin{equation}


\stdsum\int_{Q} a_{ij}(x) u_{x_j} \psi_{x_i}


dx\,dt+\int_{Q} u_t(x,t)\psi(x,t)dx\,dt+\int_{Q}


f(x,t,u)\psi(x,t) dx\,dt=0 \label{ii}


\end{equation}


и при любой функции $\chi(x,t)\in C^\infty(Q)$, равной нулю в


окрестности $\Upsilon_T(\Omega)$, выполнено включение


$\chi(x,t)(u(x,t)-\varphi(x,t))\in\stackrel{0}{W_2^{1,1}}(Q)$.


\end{defn}





Вопросы разрешимости, а также единственности решения


задачи \eqref{P}, \eqref{B} изучены в (\cite{LSU}, гл.\,V, \S\,6, с.\,524).


\medskip





{\bf 1.} Пусть в области $Q$ лежат $m$ цилиндров вида


$Q(y_{k,\varepsilon},0, R_\varepsilon,T)$ с центрами в точках


$y_{k,\eps}$ одинакового радиуса $R_{\eps}$ и высотами $T$,


$\eps=m^{-1}$ --- малый параметр, $m$ --- натуральное число;


$k=1,\dotsc,m$. Будем предполагать, что расстояние между центрами


этих шаров не меньше, чем~$4R_\eps$.


Положим $Q_\eps=Q \setminus Q(R_\eps, T)$, где $Q(R_\eps,


T)=\bigcup\limits_{k=1}^m Q(y_{k,\varepsilon},0,


R_\varepsilon,T)$.





Пусть $u_\eps(x,t)$~-- решение задачи


\begin{alignat}{2}


\cal{L}u_\eps(x,t)&=a(x)|u_\eps(x,t)|^{\sigma-1}u_\eps(x,t),


&\quad (x,t)&\in Q_\eps, \label{P2}\\


%


u_\eps(x,t)&=\varphi(x,t), &\quad (x,t)&\in S(Q), \label{B2}\\


%


u_\eps(x,0)&=u_0(x), &\quad x&\in\Omega, \label{I2}


\end{alignat}


где $a(x)$ --- ограниченная измеримая функция,


$a(x)\ge a_0=\const>0$, $\varphi(x)\in W_{2}^{1,1}(Q_\eps)$.








\begin{defn}


Функция $u_\eps(x,t)\in W^{1,1}_{2,\loc}(Q_\eps)\bigcap C(Q_\eps)$


называется {\it обобщенным решением задачи}


\eqref{P2}--\eqref{I2}, если для любой функции


$\psi(x,t)\in \overset{0}{W}{}_{2}^{1,0}(Q_\eps)$ она удовлетворяет


интегральному соотношению


$$


\stdsum\int_{Q_\eps} a_{ij}(x) u_{x_j} \psi_{x_i}


dx\,dt+\int_{Q_\eps}


u_t(x,t)\psi(x,t)dx\,dt+\int_{Q_\eps}


a(x)|u_\eps(x,t)|^{\sigma-1}u_\eps(x,t)\psi(x,t)dx\,dt=0,


$$


принимает значение $u_0(x)$ при $t=0$ и при любой функции


$\chi(x,t)\in C^\infty(Q)$, равной нулю в окрестностях


$Q(R_\eps,T)$ и $\Upsilon_T(\Omega)$, выполнено включение


$\chi(x,t)(u_\eps(x,t)-\varphi(x,t))\in\overset{0}{W}{}_2^{1,1}(Q_\eps)$.


\end{defn}





Для задачи \eqref{P2}--\eqref{I2} не выполняются, вообще


говоря, теоремы единственности, т.\,к. на границах дырок


перфорированной области $Q_\eps$ не заданы никакие граничные


условия. Покажем, что при некоторых условиях предел решений


задачи \eqref{P2}--\eqref{I2} будет существовать, причем


этот предел не зависит от выбора последовательности $u_\eps(x,t)$;


здесь $u_\eps(x,t)$ --- любое решение


задачи \eqref{P2}--\eqref{I2}.





Пусть $u(x,t)$ --- решение следующей задачи:


\begin{alignat}{2}


\cal{L}u(x,t)&=a(x)|u(x,t)|^{\sigma-1}u(x,t), &\quad (x,t)&\in Q, \label{P1}\\


u(x,t)&=\varphi(x,t), &\quad (x,t)&\in S(Q),\label{B1}\\


u(x,0)&=u_0(x), &\quad x&\in\Omega. \label{I1}


\end{alignat}








\begin{thm}


Пусть $\sigma>\frac{n}{n-2}$ и


\begin{equation}


R_\eps\sim\varepsilon^\alpha \ \text{ при } \ \eps\to 0,\quad


\text{где } \ \alpha>\frac{\sigma-1}{(n-2)\sigma-n}.


\label{Reps}


\end{equation}


Тогда решение задачи \eqref{P2}--\eqref{I2} стремится к


решению задачи \eqref{P1}--\eqref{I1} в следующем


смысле\/${:}$ для любого положительного постоянного $d$


\begin{equation}


\sup_{x\in \Omega_{\eps,d},\ t\in [0,T]}|u_\eps(x,t)-u(x,t)|\to


0\ \text{ при } \ \eps\to 0, \label{appr}


\end{equation}


где


$$


\Omega_{\eps,d}=\{x\in \Omega:\dist(x,\dd


B(y_{k,\eps},R_\eps))>d;\ k=1,\dotsc,m\}.


$$


\end{thm}





Отметим, что при этом общий объем дырок $\Vol[Q(R_\eps,T)]$


стремится к нулю при $\eps\to 0$.








\begin{cor}


Если в перфорированной области $Q_\eps$ существует множество,


отделенное от всех дырок фиксированной постоянной, то на этом


множестве сходимость будет равномерной.


\end{cor}








\begin{lem}


Для любых $a,b\in\bbR$ и $\sigma>1$ имеет место неравенство


$$


\big||a|^{\sigma-1}a-|b|^{\sigma-1}b\big|\ge


C_\sigma|a-b|^\sigma=C_\sigma|a-b|^{\sigma-1}|a-b|,


$$


где $C_\sigma$ --- некоторая положительная постоянная, зависящая


от $\sigma$, но не зависящая от $a$ и $b$. \label{L1}


\end{lem}








\begin{pf*}{Доказательство теоремы 1}


Доказательство проводится так же, как в работах \cite{MP1},


\cite{MP2} для эллиптических операторов. Для разности


$w_\eps(x,t)=u_\eps(x,t)-u(x,t)$ имеем


$$


\cal{L}w_\eps(x,t)-a(x)


(|u_\eps(x,t)|^{\sigma-1}u_\eps(x,t)-|u(x,t)|^{\sigma-1}u(x,t))=0.


$$


Умножим это равенство на $\sgn w_\eps(x,t)$ и применим лемму 1, тогда


\begin{gather*}


0=\sgn w_\eps(x,t)\cal{L}w_\eps(x,t)-a(x)


\big||u_\eps(x,t)|^{\sigma-1}u_\eps(x,t)-|u(x,t)|^{\sigma-1}u(x,t)\big|\le\\


%


\le\sgn w_\eps(x,t)\cal{L}w_\eps(x,t)-C_{\sigma}a_0|u_\eps(x,t)-u(x,t)|^\sigma=


\sgn w_\eps(x,t)\cal{L}w_\eps(x,t)-C_{\sigma}a_0|w_\eps(x,t)|^\sigma,


\quad (x,t)\in Q_\eps.


\end{gather*}


Таким образом, функция $w_\eps(x,t)$ удовлетворяет неравенству


\begin{equation}


\sgn w_\eps(x,t)


\cal{L}w_\eps(x,t)-b_0|w_\eps(x,t)|^\sigma\ge0,\ \ (x,t)\in


Q_\eps, \label{nerlk}


\end{equation}


с граничным условием


\begin{equation}


w_\eps(x,t)=0,\quad (x,t)\in


S(Q)=\partial\Omega\times[0,T],\label{wo}


\end{equation}


и с нулевым начальным условием


\begin{equation}


w_\eps(x,0)=0,\quad x\in\Omega; \label{w02}


\end{equation}


здесь $b_0=C_{\sigma}a_0$.





Согласно теореме 8.2 работы \cite{Ch} для решений полулинейных


параболических уравнений вида \eqref{P1} имеет место оценка


$$


|w_\eps(x,t)|\le C\big(\min(R^2, T)\big)^{\frac1{1-\sigma}},


$$


где $R$, $T$ --- соответственно радиус и высота цилиндра, в котором


определено решение $u(x,t)$ уравнения \eqref{P1}. Так как


коэффициенты оператора не зависят от $t$, то из последней оценки


можно исключить зависимость от $T$, поскольку $w(x,t)$ обращается


в нуль при $t=0$. Действительно, $w(x,t)$ всегда можно считать


решением в достаточно высоком цилиндре фиксированного радиуса,


если доопределить $w(x,t)$ нулем при $t<0$. Следовательно,


\begin{equation}


|w_\eps(x,t)|\le C\big(\min(R^2,


0)\big)^{\frac1{1-\sigma}}=CR^{\frac2{1-\sigma}}. \label{u}


\end{equation}





Аналогично работе \cite{KL} эту оценку можно получить и для


решений неравенств вида \eqref{nerlk}. Буквой $C$ здесь и далее


будем обозначать любую положительную постоянную, не зависящую от


$x$, $t$ и $\eps$.





Предположим, что $u_\eps(x,t)$ не стремится к $u(x,t)$ в


смысле \eqref{appr} при $\eps\to 0$. Тогда найдется положительная


постоянная $C_0$ такая, что для любого $\eps_0=m_0^{-1}$


существуют $\eps=m^{-1}<\eps_0$ и $(x_\eps,t)\in


\Omega_{\eps,d}\times [0,T]$ такие, что $|w_\eps(x_\eps,t)|>C_0$,


т.\,е. $w_\eps(x_\eps,t)>C_0$ либо $w_\eps(x_\eps,t)<-C_0$. Будем


считать, что первый случай реализуется при различных $\eps$


бесконечно много раз (для второго случая $w_\eps(x_\eps,t)<-C_0$


рассмотрение аналогично).





Обозначим через $\Gamma(x,x_0)$ фундаментальное решение


оператора ${\cal{L}}$, имеющее особенность в точке $x_0$,


т.\,е. $\cal{L}\Gamma(x)=\delta(x_0)$. Известно ~\cite{LSW},


что такое фундаментальное решение существует и удовлетворяет


неравенствам


\begin{equation}


C_1|x-x_0|^{2-n} \le\Gamma(x,x_0)\le C_2|x-x_0|^{2-n}, \label{fs}


\end{equation}


где положительные постоянные $C_1$ и $C_2$ зависят от $\lambda$ из


условия~\eqref{lmbd}.





Положим


\begin{equation}


g_\eps(x)=\frac{1}{m}\sum_{k=1}^m\Gamma(x-y_{k,\eps}).


\label{gdef}


\end{equation}





Рассмотрим функцию


\begin{equation}


v_\eps(x,t)=w_\eps(x,t)-h_{\eps}(x), \label{vdef}


\end{equation}


где $h_{\eps}(x)=\frac{C_0}{2}\frac{g_\eps(x)}{g_\eps(x_\eps)}$.


Тогда выполняется $v_\eps(x_\eps,t)=\frac{C_0}{2}>0$.





Покажем, используя принцип максимума, что это невозможно.


Согласно \eqref{u} на границах цилиндров


$Q(y_{k,\eps},0,2R_\eps,T)=B(y_{k,\eps},2R_{\eps})\times [0,T]$,


коаксиальных с


$Q(y_{k,\eps},0,R_\eps,T)=B(y_{k,\eps},R_{\eps})\times [0,T]$,


выполнено неравенство


\begin{equation}


|w_\eps(x,t)|<CR_\eps^{\frac{2}{1-\sigma}}. \label{10_o}


\end{equation}


Из нижней оценки \eqref{fs} сделаем заключение, что на границах шаров


$B(y_{k,\eps},2R_{\eps})$ имеет место


\begin{equation}


g_\eps(x)>C'_1\eps R_\eps^{2-n},


\label{10_7}


\end{equation}


где $C'_1=2^{2-n}C_1$. При этом из \eqref{Reps} следует


\begin{equation}


R_\eps^{\frac{2}{1-\sigma}}=


\bold o(\eps R_\eps^{2-n})\ \text{ при } \ \eps\to 0. \label{oo}


\end{equation}





Таким образом, получаем, что при $\sigma>\frac{n}{n-2}$ для


выполнения \eqref{oo} необходимо, чтобы\linebreak


$\alpha>\frac{\sigma-1}{(n-2)\sigma-n}$.





Так как точка $x_\varepsilon \in\Omega_{\varepsilon,d}$, то


$g_\eps(x_\eps,)<Cd^{2-n}$. Поэтому в силу \eqref{oo} на границах


цилиндров $Q(y_{k,\eps},0,2R_\eps,T)$ будет выполнено неравенство


$w_\eps(x,t)<g_\eps(x)$, т.\,е. выполнено $v_\eps(x,t)<0$.


При $(x,t)\in\dd\Omega\times [0,T]$ согласно \eqref{wo}


и \eqref{gdef} имеем $w_\eps(x,t)=0$, $g_\eps(x)\geq 0$ и


поэтому $v_\eps(x,t)\leq 0$ при $(x,t)\in S(Q)$. При $t=0$


имеем $w_\eps(x,t)=0$, $v_\eps(x,t)\leq 0$.


Кроме того, из оценок \eqref{u}, \eqref{fs} с учетом


определения \eqref{gdef} следует, что $w_\eps(x,t)\to0$ и


$g_\eps(x)\to0$ при $|x|\to\infty$. Таким образом, из принципа


максимума заключаем, что $v_\eps(x,t)\leq 0$ на параболической


границе $\Gamma(Q)$.





Определим множество $D\equiv\{(x,t) \in \Omega_{\eps}\times [0,T]:


v_\eps(x,t)>0\}\ne\emptyset$. Обозначим через $D_0$ компоненту связности


множества $D$, содержащую точку $(x_\eps,t)$.





Покажем, что $\cal{L}v_\eps(x,t)>0$ при $(x,t)\in D$.


Действительно, по определению


множества $D$ \ $u_\eps(x,t)-u(x,t)=w_\eps(x,t)>h_\eps(x)>0$,


т.\,е. $u_\eps(x,t)>u(x,t)$ при $(x,t)\in D$. Тогда


\begin{equation}


\cal{L}v_\eps(x,t)=\cal{L}(u_\eps(x,t)-u(x,t))-\cal{L}h_\eps(x)=a(x)


\big(|u_\eps(x,t)|^{\sigma-1}u_\eps(x,t)-|u(x,t)|^{\sigma-1}u(x,t)\big)>0.


\label{MaxPrince}


\end{equation}


Поскольку $v_\eps(x_\eps,t)>0$ согласно \eqref{vdef},


$v_\eps(x,t)<0$ при $x\in\dd(Q(2R_{\eps}, T))$, $v_\eps(x,0)<0$


и


$v_\eps(x,t)\to0$ при $|x|\to\infty$, то функция $v_\eps(x,t)$


достигает локального максимума во внутренней точке области $D_0$,


что с учетом~\eqref{MaxPrince} противоречит принципу максимума


для функции $v_\eps(x,t)$.





Следовательно, наше предположение неверно, т.\,е. $u_\eps(x,t)$


стремится к $u(x,t)$ в смысле \eqref{appr} при $\eps\to 0$.


\end{pf*}





Для объема шара радиуса $R$ в пространстве $\Rn$ выполнено


равенство $\Vol[B(0, R)]=\mu_nR^n$, где постоянная $\mu_n$ ---


объем единичного шара в $\bbR^n$, зависит только от $n$. Оценим


общий объем дырок:


$$


\bigg| \bigcup\limits_{k=1}^m Q(y_{k,\eps},0,R_\eps,T) \bigg| \le


m \Vol[Q(0,0,R_\eps,T)]=T\mu_nmR_{\eps}^n=\bold o


\big(\eps^{\frac{2\sigma}{(n-2)\sigma-n}}\big)


\to 0 \ \text{ при } \ \eps\to 0.


$$








{\bf 2.} Будем теперь предполагать, что расстояние между


центрами этих шаров не менее, чем $8R_\eps$.





Положим $Q(4R_\eps,T)=\bigcup\limits_{k=1}^m


Q(y_{k,\eps},0,4R_\eps, T)$.








\begin{lem}[{\series{m}\shape{n}\selectfont oбобщенное неравенство Юнга}]


Пусть $a,b,\beta>0$, $0<\alpha<1$. Тогда


$$


ab\le\frac{\alpha}{\beta}a^{\frac{1}{\alpha}}+


\frac{1-\alpha}{\beta^{-\frac{\alpha}{1-\alpha}}}b^{\frac{1}{1-\alpha}}.


$$


В частности, при $\alpha=\frac{1}{2}$ имеем


$2ab\le\frac{a^2}{\beta}+\beta b^2$. \label{L2}


\end{lem}








\begin{thm}


При любом $\sigma>1$ и $R_\eps=O(\eps^\gamma)$ при ${\eps\to0}$,


где $\gamma>\frac{1}{n-2}$, выполнено


$$


\int\limits_{Q\setminus Q(4R_\eps,T)}


\frac{|\grad(u_\eps(x,t)-u(x,t))|^2+|u_\eps(x,t)-u(x,t)|^{\sigma+1}}


{1+|u_\eps(x,t)-u(x,t)|^2}


dx\,dt\to 0\ \text{ при }\ \eps\to 0.


$$


\label{Th2}


\end{thm}





\begin{pf}


Рассмотрим функцию $\theta_\eps(x)\in C^\infty(\Omega)$ такую, что


$0\le\theta_\eps(x)\le 1$, $x\in\Omega$,


$$


\theta_\eps(x) =


\begin{cases}


1,& x\in \Omega\setminus \mathop{\cup}\limits_{k=1}^m B(y_{k,\varepsilon},4R_\eps); \\


0,& x\in \mathop{\cup}\limits_{k=1}^m B(y_{k,\varepsilon},2R_\eps),


\end{cases}


$$


$|\grad\theta_\eps(x)|<\frac{C_\Omega}{R_\eps}$, $x\in\Omega$;


здесь $C_\Omega$ зависит от области $\Omega$ и не зависит от $\eps$.





Рассмотрим функцию


\begin{equation}


\psi(x,t)=\theta_\eps^p(x)\eta(u_\eps(x,t)-u(x,t)), \label{psi}


\end{equation}


где


$\eta(\tau)=


\begin{cases}


\tau, &|\tau|\le 1;\\


|\tau|^{q-1}\tau, &|\tau|>1,


\end{cases}


$ \ $q>-1$ таков, что $\sigma>\frac{n+2q}{n-2}$;


отметим, что такое $q$ можно подобрать


при любом $\sigma>1$, $p=p(q)>2$ --- решение 


уравнения $p=\frac{(p-2)(\sigma+q)}{1+q}$, т.\,е.


\begin{equation}


p=2\frac{\sigma+q}{\sigma-1}, \label{8}


\end{equation}


требуемое неравенство $p>2$ выполнено при требуемом условии на $q>-1$.


Отметим также, что функцию $\psi(x,t)$ можно подставлять в


интегральное тождество из определения обобщенного решения, т.\,к.


она принадлежит классу $W^{1,0}_{2}(Q_{\eps})$ в силу правила


цепного дифференцирования (\cite{GT}, с.\,151) и обращается в нуль на границе


области $Q_{\eps}$, на границе области $Q$ за счет того, что


разница $u_\eps(x,t)-u(x,t)$ обращается в нуль на $\partial Q$, а


на границе дырок за счет множителя $\theta_\eps^p(x)$. Отметим


также, что поскольку на всех дырках этот множитель равен нулю, в


интегральном тождестве мы можем в качестве области интегрирования


писать $Q$.





Как и в доказательстве теоремы 1, обозначим


$w_\eps(x,t)=u_\eps(x,t)-u(x,t)$. Для уравнения


$$


Lw_\eps(x,t)-\frac {\dd w_\eps(x,t)}{\dd t}=a(x)


(|u_\eps(x,t)|^{\sigma-1}u_\eps(x,t)-|u(x,t)|^{\sigma-1}u(x,t))


$$


запишем интегральное тождество, взяв в качестве пробной функции


$\psi(x,t)$ из \eqref{psi}, получим


\begin{multline}


\int_{Q}\bigg(\stdsum\frac{\dd}{\dd x_i}\bigg(a_{ij}(x)


\frac{\dd w_\eps(x,t)}{\dd x_j}\bigg)-\frac {\dd


w_\eps(x,t)}{\dd t}\bigg)\psi(x,t)dx\,dt=\\


=\int_{Q}a(x)\big(|u_\eps(x,t)|^{\sigma-1}


u_\eps(x,t)-|u(x,t)|^{\sigma-1}u(x,t)\big) \psi(x,t) dx\,dt.\label{9}


\end{multline}





Обозначим $Q_1=\{(x,t)\in Q:|u_\eps(x,t)-u(x,t)|\le 1 \}\subset Q$


и


$Q_2=\{(x,t)\in Q:\linebreak |u_\eps(x,t)-u(x,t)|> 1 \}\subset Q$.





Оценим правую часть равенства \eqref{9}, используя лемму 1 при


$a=u_\eps(x,t)$, $b=u(x,t)$,


\begin{multline}


\label{10}


\int_{Q}a(x)\big(|u_\eps(x,t)|^{\sigma-1}u_\eps(x,t)-


|u(x,t)|^{\sigma-1}u(x,t)\big)\psi(x,t) dx\,dt =\\


=\int_{Q_1}a(x)(|u_\eps(x,t)|^{\sigma-1}u_\eps(x,t)-|u(x,t)|^{\sigma-1}u(x,t))


w_\eps(x,t)\tmp dx\,dt + \\


+\int_{Q_2}a(x)(|u_\eps(x,t)|^{\sigma-1}u_\eps(x,t)-u(x,t)^{\sigma-1}u(x,t))


|w_\eps(x,t)|^{q-1}w_\eps(x,t)\tmp dx\,dt \ge \\ 


\ge a_0C_\sigma\int_{Q_1}|w_\eps(x,t)|^{\sigma+1}\tmp dx\,dt


+a_0C_\sigma\int_{Q_2} |w_\eps(x,t)|^{\sigma+q}\tmp dx\,dt.


\end{multline}


После тождественных преобразований и интегрирования по частям в


левой части равенства~\eqref{9} получим


\begin{align}


J_1+J_2+J_3+J_4 \equiv &-\int_{Q_1}\stdsum a_{ij}(x)\frac


{\dd w_\eps(x,t)}{\dd x_j}\frac{\dd w_\eps(x,t)}{\dd x_i}\tmp dx\,dt -\notag \\


&-\int_{Q_1}\stdsum a_{ij}(x)\frac {\dd w_\eps(x,t)}{\dd


x_j}w_\eps(x,t)p\theta_\eps^{p-1}(x)\frac{\dd \tm}{\dd x_i} dx\,dt -\notag \\


&-\int_{Q_2}\stdsum a_{ij}(x) \frac {\dd w_\eps(x,t)}{\dd


x_j}q|w_\eps(x,t)|^{q-1}\frac {\dd w_\eps(x,t)}{\dd x_i}\tmp dx\,dt - \notag\\


\label{11}


&-\int_{Q_2}\stdsum a_{ij}(x) \frac {\dd w_\eps(x,t)}{\dd


x_j}|w_\eps(x,t)|^{q-1}w_\eps(x,t)p\tmd{p-1}\frac{\dd \tm}{\dd


x_i} dx\,dt. 


\end{align}


Оценим первое и третье слагаемые правой части


тождества~\eqref{11}, учитывая~\eqref{lmbd},


\begin{align}


J_1&\le -\frac{1}{\lambda}\int_{Q_1}\mgrad^2\tmp dx\,dt; \label{12} \\


%


J_3 &\le -\frac{q}{\lambda}\int_{Q_2}\mgrad^2


|w_\eps(x,t)|^{q-1}\tmp dx. \label{13}


\end{align}


Для оценки второго и четвертого слагаемых правой части


тождества \eqref{11} используется неравенство Юнга при


$\alpha=\frac{1}{2}$, $a=\mgrad$, $b=|\grad


\tm|$. Выбирая $\beta=\frac{p\lambda^2}{\tm}$ и


$\beta=\frac{p\lambda^2|w_\eps(x,t)|}{q\tm}$ при каждом


$(x,t)$, соответственно получаем


\begin{align}


J_2 &\le p\lambda \int_{Q_1} \mgrad \tgrad \tmd{p-1} dx\,dt\le \notag\\


%


& \le p\lambda \int_{Q_1} \tmd{p-1}


\bigg(\frac{\tm}{2p\lambda^2}\mgrad^2


+\frac{p\lambda^2}{2\tm}\tgrad^2 \bigg) dx\,dt= \notag\\


&=\frac{1}{2\lambda} \int_{Q_1} \mgrad^2\tmp dxdt +C_4


\int_{Q_1} \tgrad^2 \tmd{p-2} dx\,dt; \label{14} \\


%


J_4 &\le p\lambda \int_{Q_2}\mgrad \tgrad |w_\eps(x,t)|^q \tmd{p-1}\le \notag\\


&\le p\lambda \int_{Q_2} \tmd{p-1}|w_\eps(x,t)|^q\bigg(


\frac{q\tm}{2p\lambda^2|w_\eps(x,t)|}\mgrad^2


+\frac{p\lambda^2|w_\eps(x,t)|}{2q\tm}\tgrad^2\bigg)dx\,dt = \notag\\


%


&=\frac{q}{2\lambda}\int_{Q_2}


\mgrad^2|w_\eps(x,t)|^{q-1}\tmp dx +C_5


\int_{Q_2}\tgrad^2|w_\eps(x,t)|^{q+1}\tmd{p-2} dx\,dt,


\label{15}


\end{align}


где $C_4=\frac{p^2\lambda^3}2$, $C_5=\frac{p^2\lambda^3}{2q}$.





Теперь оценим второе слагаемое правой части


неравенства \eqref{15}. Полагая в неравенстве Юнга


$a=|w_\eps(x,t)|^{q+1}\tmd{p-2}$, $b=\tgrad^2$ при


$\alpha=\frac{1+q}{\sigma+q}$ и выбрав


$\beta=\frac{2C_5}{a_0C_\sigma}\alpha$, с учетом~\eqref{8} получаем


\begin{equation}


C_5 \int_{Q_2} \tgrad^2 |w_\eps(x,t)|^{q+1} \tmd{p-2} dx\,dt \le 


%


\frac{a_0C_\sigma}{2}\int_{Q_2}|w_\eps(x,t)|^{\sigma+q}\tmp


dx\,dt+ C_6\int_{Q_2}\tgrad^{p} dx\,dt, \label{16}


\end{equation}


где $C_6=C_5\frac{\sigma-1}{\sigma+q}


\big(\frac{2C_5}{a_0C_\sigma}\frac{1+q}{\sigma+q}\big)^


{\frac{q+1}{\sigma-1}}=\frac{\sigma-1}{\sigma+q}


\big(\frac{2(1+q)}{a_0C_\sigma(\sigma+q)}\big)^


{\frac{q+1}{\sigma-1}}C_5^{\frac{\sigma+q}{\sigma-1}}$.





Рассмотрим теперь член, содержащий производные по $t$. Для этого


заметим, что функция $\eta(\tau)$, входящая в определение


$\psi(x,t)$, обладает первообразной


$$


\eta_1(\tau)=


\begin{cases}


\frac{1}{2}\tau^2,& |\tau| \le 1; \\


\frac{1}{q+1}|\tau|^{q+1},& |\tau| > 1.\\


\end{cases}


$$


Действительно,


$$


\bigg(\frac{1}{q+1}|\tau|^{q+1}\bigg)_\tau=


\frac{1}{2}\frac{2}{q+1}(|\tau|^{2\frac{q+1}{2}})_\tau


=\frac{1}{2}(|\tau|^2)_\tau(|\tau|^2)^{\frac{q-1}{2}}=


\frac{1}{2}\tau(|\tau|^2)^{\frac{q-1}{2}}=\eta(\tau).


$$


Поэтому член, содержащий производные по $t$, можем записать в виде


\begin{multline}


-\int_{Q}\frac{\dd w_\eps(x,t)}{\dd t}\psi(x,t) dx\,dt


=-\int_{Q}\frac{\dd w_\eps(x,t)}{\dd t}\eta(w_\eps(x,t))\tmp dx\,dt=\\


=-\int_{Q}\frac{\dd}{\dd t} ( \eta_1(w_\eps(x,t))


\tmp) dx\,dt = -\int_{\Omega}(


\eta_1(w_\eps(x,t))\tmp) \big|^{T}_{0} dx\leq 0.


\label{pttt}


\end{multline}


Последнее выполняется, поскольку $\eta_1(x,t)\tmp$ положительно и


обращается в нуль при $t=0$.


Учитывая равенства \eqref{9}, \eqref{11} и


оценки \eqref{12}--\eqref{pttt}, имеем


\begin{multline}


\int_{Q}\bigg(Lw_\eps(x,t)\psi(x,t)-\frac{\dd


w_\eps(x,t)}{\dd t}\psi(x,t)\bigg) dx\,dt \le


\int_{Q}Lw_\eps(x,t)\psi(x,t) dx\,dt \le \\


%


\le -\frac{1}{2\lambda}\int_{Q_1}\mgrad^2\tmp dx\,dt-


\frac{q}{2\lambda}\int_{Q_2}


\mgrad^2|w_\eps(x,t)|^{q-1}\tmp dx\,dt +\\


%


+C_4 \int_{Q_1}\tgrad^2 \theta_\eps^{p-2}(x) dx\,dt +


\frac{a_0C_\sigma}{2}\int_{Q_2}|w_\eps(x,t)|^{\sigma+q}\tmp


dx\,dt+ C_6\int_{Q_2}\tgrad^{p} dx\,dt. \label{17}


\end{multline}


Объединяя \eqref{9}, \eqref{10} и \eqref{17}, получаем


\begin{multline*}


-\frac{1}{2\lambda}\int_{Q_1}\mgrad^2\tmp dx\,dt -


\frac{q}{2\lambda}\int_{Q_2} \mgrad^2


|w_\eps(x,t)|^{q-1}\tmp dx\,dt +\\


%


+C_4 \int_{Q_1}\tgrad^2 \theta_\eps^{p-2}(x) dx\,dt +


\frac{a_0C_\sigma}{2}\int_{Q_2}


|w_\eps(x,t)|^{\sigma+q}\tmp dx\,dt+


C_6\int_{Q_2}\tgrad^{2\frac{\sigma+q}{\sigma-1}} dx\,dt \ge \\


%


\ge a_0C_\sigma\int_{Q_1} |w_\eps(x,t)|^{\sigma+1}\tmp


dx\,dt + a_0C_\sigma\int_{Q_2}|w_\eps(x,t)|^{\sigma+q}\tmp dx\,dt.


\end{multline*}


После элементарных преобразований получаем


\begin{multline}


C_4 \int_{Q_1}\tgrad^2 \theta_\eps^{p-2}(x) dx\,dt +


C_6 \int_{Q_2}\tgrad^{p} dx\,dt \ge \\


%


\ge a_0C_\sigma \int_{Q_1} |w_\eps(x,t)|^{\sigma+1}\tmp


dx\,dt +\frac{1}{2\lambda} \int_{Q_2}\mgrad^2\tmp dx\,dt + \\


%


+\frac{a_0C_\sigma}{2} \int_{Q_2}


|w_\eps(x,t)|^{\sigma+q}\tmp dx\,dt + \frac{q}{2\lambda}


\int_{Q_2} \mgrad^2|w_\eps(x,t)|^{q-1}\tmp dx\,dt \ge \\


%


\ge C_7 \int_{Q} \frac{|w_\eps(x,t)|^{\sigma+1}}


{1+|w_\eps(x,t)|^2}\tmp dx\,dt + C_8


\int_{Q}\frac{\mgrad^2}{1+|w_\eps(x,t)|^2}\tmp dx\,dt \ge \\


%


\ge C_9 \int\limits_{Q \setminus Q(4R_\eps,T)} \frac{|\grad


w_\eps(x,t)|^2+|w_\eps(x,t)|^{\sigma+1}} {1+|w_\eps(x,t)|^2} dx\,dt,


\label{18}


\end{multline}


где $C_7=\frac{a_0C_\sigma}2$, $C_8=\min\big\{\frac1{2\lambda},\frac


q{2\lambda}\big\}$, $C_9=\min\{C_7, C_8\}$.





Предпоследнее неравенство получено из следующего утверждения.


Если $0\leq a\leq 1$, то выполнено неравенство $1 \geq


\frac{1}{1+a^2}$; если же $ a>1$, то выполнено


$a^{q-1}>\frac{1}{a^2}>\frac{1}{1+a^2}$ при $q>-1$.





Теперь докажем, что


\begin{align}


\int_{Q_1}\tgrad^2 \theta_\eps^{p-2}(x)


dx\,dt&\to0 \ \text{ при }\ \eps\to 0 \label{19} \\


%


\intertext{и}


%


\int_{Q_2}\tgrad^{p}dx\,dt&\to 0 \ \text{ при }


\ \eps\to0.\label{20}


\end{align}


Действительно, т.\,к. $p>2$, то


$$


\int_{Q_1}\tgrad^2 \theta_\eps^{p-2}(x) dx\,dt


\le\int_{Q}\tgrad^2 dx\,dt \le C_{10}mR_\eps^n


\bigg(\frac{1}{R_\eps}\bigg)^2 =


O(\eps^{-1+{(n-2)\gamma}}) \to 0


$$


при $\eps\to 0$, если только $\gamma>\frac{1}{n-2}$, где


$C_{10}=TC_\Omega^2 \mu_n 2^n$;


$$


\int_{Q_2}\tgrad^{p} dx\,dt \le\int_{Q}\tgrad^{p}


dx\,dt \le C_{11}mR_{\eps}^n \bigg(\frac{1}{R_\eps}\bigg)^{p}=


C_{11}\eps^{-1}R_\eps^{n-p}= O(\eps^{-1+(n-p)\gamma})\to 0


$$


при $\eps\to 0$, т.\,к. $\gamma>\frac{1}{n-p}>\frac{1}{n-2}$;


здесь $C_{11}=TC_\Omega^{2\frac{\sigma+q}{\sigma-1}}\mu_n2^n$,


$\mu_n$ --- объем единичного $n$-мерного шара.





Из \eqref{18}, \eqref{19} и \eqref{20} следует справедливость


теоремы 2.


\end{pf}











\begin{thebibliography}{99}





\bibitem{BLP}


Bensoussan A., Lions J.-L., Papanicolaou G. {\it Asymtotic analysis for


periodic structures}. -- Amsterdam: North-Holland, 1978. -- 700~p.


\bibitem{ZhKO}


Жиков В.В., Козлов С.М., Олейник О.А. {\it Усреднение дифференциальных


операторов}. -- М.: Физматлит, 1993. -- 464~c.


\bibitem{MKh}


Марченко В.А., Хруслов Е.Я. {\it Краевые задачи в областях с мелкозернистой


границей}. -- Киев: Наук. думка, 1974. -- 280~c.


\bibitem{BF}


Brezis H., Friedman A. {\it Nonlinear parabolic equation


involving measures as initial conditions} // J. Math. Pures Appl.


-- 1983. -- V.\,62. -- \No\,1. -- P.\,73--97.


\bibitem{BPT}


Brezis H., Peletier L.A., Terman D. {\it A very singular solutions


of the heat equation with absorption} // Arch. Ration. Mech. Anal.


-- 1986. -- V.\,95. -- \No\,3. -- P.\,185--209.


\bibitem{KP}


Kamin S., Peletier L.A. {\it Singular solutions of the heat equation


with absorption} // Proc. Amer. Math. Society. -- 1985. -- V.\,95. -- \No\,2.


-- P.\,205--210.


\bibitem{IOK}


Ильин А.М., Олейник О.А., Калашников А.С. {\it Линейные уравнения


второго порядка параболического типа} // УМН. -- 1962. -- Т.\,17. -- \No\,3.


-- С.\,3--146.


\bibitem{MP1}


Матевосян О.А., Пикулин С.В. {\it Об усреднении слабонелинейных


дивергентных эллиптических операторов в перфорированном кубе}


// Матем. заметки. -- 2000. -- Т.\,68. -- \No\,3. -- С.\,390--398.


\bibitem{MP2}


Матевосян О.А., Пикулин С.В. {\it Об усреднении полулинейных


эллиптических операторов в перфорированных областях} // Матем.


сб. -- 2002. -- Т.\,193. -- \No\,3. -- С.\,101--114.


\bibitem{KL}


Кондратьев В.А., Ландис Е.М. {\it О качественных


свойствах решений одного нелинейного уравнения второго порядка} //


Матем. сб. -- 1988. -- Т.\,135. -- \No\,3. -- С.\,346--360.


\bibitem{Ch}


Чистяков В.В. {\it О свойствах решений полулинейных


параболических уравнений второго порядка} // Тр. семинара


И.Г.\,Петровского. -- 1991. -- Вып.\,15. -- С.\,70-107.


\bibitem{LSU}


Ладыженская О.А., Солонников В.А., Уральцева


Н.Н. {\it Линейные и квазилинейные уравнения параболического


типа}. -- М.: Наука, 1967. -- 736 c.


\bibitem{LSW}


Littman W., Stampacchia G., Weinberger B. {\it Regular points


for elliptic equations with discontinuous coefficients} // Ann.


Scuola Norm. Super. Pisa. Ser.\,3. -- 1963. -- V.\,17. -- \No\,1--2.


-- P.\,43--77.


\bibitem{GT}


Гилбарг Д., Трудингер Н. {\it Эллиптические дифференциальные


уравнения с частными производными второго порядка}. -- М.:


Наука, 1989. -- 464~c.





\end{thebibliography}





\vskip 2cm





\post{\it Московский государственный}{\it Поступила}


\post{\it университет}{30.06.2003}





\label{end}


\end{document}


